% Figaro Paper F2: The Wallet Without a Polity
% Displaced and Stateless Subjects under a Bonded Settlement Primitive
% Audience: humanitarian-economics researchers, refugee-policy scholars,
% legal philosophers of personhood, NGO policy staff working on stateless
% populations and displaced economies.
% Companion to Paper F1 (The Wallet as Legal Subject), which treats the
% same wallet-collapse argument from the labor-law side.

\documentclass[11pt,a4paper]{article}

\usepackage[utf8]{inputenc}
\usepackage[T1]{fontenc}
\usepackage{amsmath,amssymb,amsthm}
\usepackage{mathtools}
\usepackage{booktabs}
\usepackage{graphicx}
\usepackage{hyperref}
\usepackage{cleveref}
\usepackage{enumitem}
\usepackage{xcolor}
\usepackage[margin=1in]{geometry}
\usepackage{natbib}
\bibliographystyle{plainnat}

\theoremstyle{remark}
\newtheorem{remark}{Remark}[section]

\newcommand{\figaro}{\textsc{Figaro}}

\title{%
  The Wallet Without a Polity:\\
  Displaced and Stateless Subjects under\\
  a Bonded Settlement Primitive%
}

\author{
  Alessandro Daliana\thanks{Corresponding author. \texttt{adaliana@figaro.org}}
}

\date{April 2026}

\begin{document}

\maketitle

% ════════════════════════════════════════════════════════════════════
\begin{abstract}
The bonded settlement primitive of the mechanism paper,
implemented and verified in the implementation paper, exhibits a
property whose consequences are most consequential not for
populations conventionally treated by labor or commercial law, but
for those who occupy the negative space of those frames: refugees,
stateless persons, undocumented migrants, and populations whose
residency is beyond the reach of stable banking and enforcement
infrastructure. The primitive's precondition is a cryptographic
key, not a recognized civil-legal subjecthood. This paper
develops the consequence.

The wallet-as-legal-subject argument extends here in a different
direction. In the labor setting the wallet collapses the Roman
\emph{res}/\emph{persona} distinction so that the productive asset
is simultaneously property and party; in the present setting it
allows the productive asset to act as party where civil-legal
subjecthood is itself absent. We do not claim that the primitive
resolves statelessness. \citet{arendt1951origins} named the right
to have rights as the fundamental human right --- a right whose
denial is a civilizational failure rather than an individual
misfortune --- and the technological infrastructure considered
here cannot grant it. What it grants is narrower: a capacity to
have commerce, where the capacity to have rights is denied. The
two should not be confused. For populations for whom the first is
denied, the second is non-trivial.

The argument is positional rather than programmatic. Real
constraints --- key custody, counterparty willingness, fiat
on-ramp, collateral access --- are named and not dismissed. The
analytical claim is that the primitive moves the
humanitarian-economics gap from an architectural impossibility to
an operational engineering question, and that this is a
consequential move worth scholarship.
\end{abstract}

\medskip
\noindent\textbf{Keywords:}
statelessness, refugee economies, displaced populations, the
wallet as legal subject, Arendt, humanitarian economics,
commerce-without-jurisdiction

% ════════════════════════════════════════════════════════════════════
\section{Introduction}
\label{sec:intro}

The legal infrastructure of commerce presumes a recognized civil
subject. Contract law requires parties whose standing the forum
will accept; banking access requires identity documents the
domestic regime issues; commercial-dispute adjudication requires
a court whose jurisdiction reaches both sides. Each of these
presumptions is so deeply embedded in the architecture of
contemporary trade that its failure mode --- populations for whom
the presumptions do not hold --- is treated by the architecture as
exception rather than as a system property.

The United Nations Refugee Agency estimated 120 million forcibly
displaced persons worldwide at the end of 2023, and approximately
4.4 million persons registered as stateless, with a substantially
larger estimated population whose statelessness is undocumented
\citep{unhcr2024globaltrends}. An overlapping and considerably
larger population lacks stable banking and enforcement access for
reasons of informal migration status, residence without papers,
or residency in jurisdictions whose financial infrastructure does
not reach them. Across these categories, the productive capacity
of well over a hundred million working-age adults is
systematically underused because the legal scaffold for engaging
it is absent. The humanitarian-economics literature on refugee
economies~\citep{betts_collier_2017} and the policy frontier
between camp and labor market~\citep{pritchett2006people}
documents the cost of this gap and the existing repertoire of
responses.

This paper argues that the bonded settlement primitive
of the mechanism paper, implemented and verified
in the implementation paper, changes the architectural posture of
the gap. The primitive operates on parties whose existence as
parties does not require civil-legal recognition. A wallet can
sign; a bonded commitment can settle; an exchange can complete ---
without a host-state court, a domestic bank, or a permission-to-work
visa sitting between the parties. The labor-law consequences of
the primitive are addressed elsewhere; the present paper develops the
consequence for populations to whom the labor-law debate is a
second-order question, because the first-order question is
whether they are recognized parties to any contract in any forum
at all.

\paragraph{What this paper is not.}
This paper does not claim that the primitive resolves
statelessness. It does not confer citizenship, restore legal
recognition, or create enforcement against a state that denies a
person's standing. It is not a substitute for the political work
that addresses those gaps --- work that remains necessary, that
this paper does not duplicate, and that no technological
infrastructure can perform. It does not predict adoption rates,
endorse any particular humanitarian program, or assess the
political consequences of widespread wallet adoption among
displaced populations. Its argument is structural: the
architectural posture of the commerce-access gap changes once a
settlement primitive whose precondition is a cryptographic key
exists.

\paragraph{Paper organization.}
\Cref{sec:primitive} summarizes the settlement primitive in the
terms relevant here. \Cref{sec:wallet-brief} restates the
wallet-as-legal-subject argument from the labor-law paper in
the form needed for the present paper. \Cref{sec:capacity}
develops the central claim across four moves: scale of the
affected population, the architecture they face, what the
primitive does and does not do, and the Arendtian distinction
between the right to have rights and the capacity to have
commerce. \Cref{sec:constraints} names the operational constraints
the primitive does not dissolve. \Cref{sec:conclusion} concludes.

% ════════════════════════════════════════════════════════════════════
\section{The Settlement Primitive}
\label{sec:primitive}

We summarize only the features relevant to the present argument.
The mechanism is established in the mechanism paper; the
Solidity implementation and formal verification
in the implementation paper.

The primitive composes two mechanisms. \emph{Asymmetric bonding}
makes defection more expensive than cooperation at every bonded
edge. \emph{Buyer dominance with atomic resolution} settles every
active order in a process simultaneously, at the buyer's
signature, with no third-party adjudicator. The first mechanism
removes the trusted enforcer; the second removes the trusted
clearing house. For the present paper, both are load-bearing:
the displaced or stateless subject lacks access to either, and
the primitive's contribution is to make both unnecessary.

A transaction between two parties is consummated by dual-signed
EIP-712 commitments. Each party signs from their own wallet. Each
party locks collateral against their own performance --- the buyer
locks $2P$, the seller locks $2G$ where $G$ is the cumulative
upstream value bonded in the process. Only the buyer can trigger
resolution. Defection is structurally more expensive than
cooperation. There is no platform between the parties, no
arbitrator authorized to resolve disputes on-chain, and no admin
key permitting external override.

The features that matter for the present paper are three. First,
the kernel does not require either party to be a recognized civil
subject; it requires them to control a key that can sign. Second,
the kernel does not require a forum to enforce the agreement; the
collateral structure is the enforcement mechanism. Third, the
kernel does not route value through a banking intermediary;
settlement is direct, wallet to wallet, atomic on resolution.

% ════════════════════════════════════════════════════════════════════
\section{The Wallet as Legal Subject (Brief)}
\label{sec:wallet-brief}

Roman private law distinguished \emph{res} (things, property,
objects that can be owned and alienated) from \emph{persona}
(persons, parties, subjects capable of acting legally). The
distinction survived and structures modern private law. The
bonded primitive collapses the
\emph{persona}/\emph{res} partition: a wallet is simultaneously a
property (a key-pair owned and transferable) and a party (a
signer, bonder, settlement-receiver). The collapse is operational
rather than metaphorical; it follows from the architecture of
dual-signed bonded commitments.

For the present paper the relevant consequence is narrower than
the labor-law one but harder. In the labor setting the
\emph{persona}/\emph{res} collapse dissolves the subordination
gradient between employer and employee, because each party signs
for itself and bonds for its own performance. In the present
setting --- where one or both parties may lack stable
\emph{persona} status in the host jurisdiction at all --- the
collapse means something stronger: the wallet provides a
party-to-the-transaction whose existence does not depend on the
prior existence of a recognized civil subject. The labor-law
argument shows that the wallet can be both property and party.
The argument here is that the wallet can be a party where no
civil-legal party exists.

This is not a claim that the wallet \emph{is} a civil-legal
subject. Civil-legal subjecthood is conferred by polities, and
its absence is a political condition. The claim is narrower: a
party to a bonded commitment can transact, settle, and accumulate
verifiable performance history without first acquiring civil-legal
subjecthood, because the kernel does not consult the polity for
permission to settle a bond.

% ════════════════════════════════════════════════════════════════════
\section{Self-Sovereignty: The Cryptographic Key as Precondition}
\label{sec:self-sovereignty}

The previous section established that the wallet can be a party
where no civil-legal party exists. The deeper observation, which
the self-sovereign-identity literature has been developing
independently of the bonded-primitive context, is that this
arrangement is a particular kind of agency that deserves a name
distinct from both \emph{civil-legal personhood} and the older
\emph{natural personhood} the philosophical tradition has
analyzed under that heading. The name we adopt, following
\citet{allen2016path} and the broader self-sovereign-identity
literature \citep{tobin_reed_2017, w3c_did_2022}, is
\emph{self-sovereignty}: the condition of an agent whose
capacity to bind, transact, and accumulate verifiable history
derives from cryptographic self-control rather than from
recognition by a polity.

\paragraph{Three acts, three preconditions discharged.}
A bonded transaction reduces to three acts at the participant's
side, and each act has historically required a civil-legal
precondition that the bonded primitive replaces with a
cryptographic one.

\emph{Signing.} The act of binding oneself to terms has
historically required a recognized capacity to enter contracts ---
adulthood, sound mind, jurisdictional standing, and the
authority to act for oneself in the relevant forum. Under the
bonded primitive the capacity to sign is the capacity to control
a private key. EIP-712 typed-structured-data signatures verify
through public-key cryptography against the address that signed;
the verification does not consult any registry of recognized
subjects. A cryptographic key is the precondition; a passport is
not.

\emph{Bonding.} The act of pledging assurance has historically
required a counterparty willing to extend credit, a court willing
to enforce a guarantee, or a recognized residency status enabling
a banking relationship. Under the bonded primitive the act of
pledging assurance is the act of locking collateral in a
permissionless token. The token's transferability is enforced by
the chain's consensus, not by any host institution. The bond's
return on cooperation is enforced by the kernel, not by a court.
A wallet that holds the bond-denomination currency can pledge
without intermediation.

\emph{Settling.} The act of receiving consideration has
historically required a banking relationship with an institution
that can credit the receiving party's account in the receiving
jurisdiction's currency. Under the bonded primitive settlement
is direct wallet-to-wallet on resolution. The chain is the
clearing layer at the architectural level. The architectural
property is unconditional; the operational property in any given
denomination is not, since the bond-denomination currency is
itself a jurisdictional artifact (USDC is freezable by Circle on
Treasury request; ETH is acquired through KYC-gated exchanges).
The bond-currency embedding constraint is developed in
\Cref{sec:constraints} below; here we register that the
architectural elimination of the banking precondition is real
but its practical realization for any specific participant
remains conditional on the bond-currency layer's own
jurisdictional reach.

The three acts together constitute a complete commercial
participation. The three preconditions historically required for
the acts --- recognized capacity, intermediated assurance,
banked settlement --- are each replaced by cryptographic
substitutes that do not require recognition by the civil-legal
order. This is what self-sovereignty in our usage means: the
agent's commercial agency is constituted by what the agent can
prove cryptographically rather than by what the agent has been
recognized as politically.

\paragraph{What self-sovereignty does and does not mean.}
The term has been used in several adjacent senses, and we owe
the reader a precise scope.

The \emph{self-sovereign-identity} (SSI) literature
\citep{allen2016path, tobin_reed_2017} is grounded in identity
specifically: the participant holds verifiable credentials about
themselves rather than depending on identity providers to attest
to their identity. \citet{allen2016path} articulates ten
principles for self-sovereign identity --- existence, control,
access, transparency, persistence, portability, interoperability,
consent, minimization, and protection --- several of which
extend beyond credential-holding into the agency-of-interaction
register: control over how identity data is used, persistence
across institutional change, portability across providers. The
``operational self-sovereignty'' sense we develop in this paper
is the commerce-specific narrowing of Allen's broader
agency-of-interaction principles, applied to the bonding,
signing, and settling acts of bilateral commerce rather than to
identity per se. The W3C Decentralized Identifiers specification
\citep{w3c_did_2022} formalizes the architectural shape:
identifiers controlled by the subject through cryptographic
keys, resolvable to documents the subject publishes, with no
central registrar. SSI in the W3C-DID register is about
\emph{identity}, and the bonded primitive does not implement
identity (the kernel knows wallets, not persons; identity is a
runtime-tier concern composed by parties who want it).

The \emph{libertarian / Lockean} sense of sovereignty
\citep{nozick1974anarchy, locke1689treatises} is normative: the
agent has natural rights of self-ownership that political
arrangements must respect or violate. The bonded primitive is
silent on natural rights; it provides architecture, not ethics.
A reader who takes self-sovereignty in this sense should read
our usage as compatible with but not requiring the normative
framework: the architecture supplies a substrate on which a
self-ownership reading \emph{can} be expressed, but the
substrate makes no claim about whether self-ownership is
morally fundamental.

The sense we adopt in this paper is narrower than either: an
agent whose commercial agency is operationally constituted by
cryptographic self-control rather than by polity recognition.
This is a structural fact about how the bonded primitive
operates, not a normative claim about how it should. The
displaced participant who sign-bonds-settles a process is
exercising self-sovereignty in this operational sense, regardless
of how their natural-rights status is read by any external
philosophical framework. The architectural fact is what we
claim; the normative reading is an interpretation that readers
in different traditions will compose differently.

\paragraph{Why this matters specifically for the displaced.}
The general operational sense of self-sovereignty applies
universally: any wallet that signs-bonds-settles is exercising
this kind of agency. But the marginal value of the architecture
is asymmetric. A participant whose civil-legal subjecthood is
already conferred and stable has plenty of paths through which
to express their commercial agency; the bonded primitive is one
substrate among many, and self-sovereignty in the operational
sense adds modestly to a position that was already operational.
A participant whose civil-legal subjecthood is denied or
unstable does not have those paths. For them, the operational
self-sovereignty the bonded primitive supplies \emph{is} the
agency, not an augmentation to it. The displaced or stateless
participant whose signing-bonding-settling is enabled by
nothing more than control of a private key has acquired a
commercial-participation capacity that the surrounding
infrastructure does not grant.

This is what the present paper has been calling the capacity to
have commerce. Self-sovereignty is the precise architectural
property that makes the capacity available to a population for
which it was not previously available. The Arendtian distinction
of the previous section --- right to have rights versus capacity
to have commerce --- decomposes once self-sovereignty is named
explicitly. The right to have rights is conferred by a polity
recognizing the agent. Self-sovereignty is constituted by the
agent's cryptographic control of their own commercial agency.
The two are different things, and conflating them has been a
recurring error in the literature this paper has been pushing
back against in both directions: those who claim cryptographic
infrastructure restores Arendtian rights overstate; those who
claim cryptographic infrastructure cannot do meaningful work for
the displaced because it does not restore Arendtian rights also
overstate. Self-sovereignty is real, modest, and operationally
load-bearing where conventional infrastructure fails; it is not
a substitute for political belonging where conventional
infrastructure succeeds.

\paragraph{Key custody as the load-bearing operational variable.}
Self-sovereignty's architectural cleanness should not obscure
its operational fragility. The architecture turns on the
participant's ability to hold and use a private key reliably.
For the displaced and stateless populations of greatest interest
to this paper, key custody is exactly where the architecture
meets the hardest constraints: unstable access to power and
connectivity, devices vulnerable to seizure or loss, environments
where social-recovery networks are themselves disrupted by the
displacement. Self-sovereignty in the operational sense
collapses to the participant's actual capacity to retain
exclusive control of their key over time. The
\Cref{sec:constraints} treatment of these constraints is
therefore not parenthetical; it is where the architectural
property cashes out into deliverable agency.

The self-sovereign-identity community has invested significant
effort in mitigations: social-recovery wallets where
participant-trusted contacts can collectively recover a lost key
\citep{buterin2021socialrecovery, argent_social_recovery,
safe_social_recovery}, secure
elements in inexpensive hardware, mnemonic backup standards
encoded in widely-spoken languages, and emerging
biometric-binding schemes that tie key control to a physical
attribute. None of these is perfect; each is an engineering
artifact with its own threat model. We do not endorse any
specific mitigation here. We name the constraint because
self-sovereignty without durable key custody is structurally
incomplete, and the engineering of durable key custody for
displaced populations is the work that translates the
architectural availability into deliverable agency.

% ════════════════════════════════════════════════════════════════════
\section{The Capacity to Have Commerce}
\label{sec:capacity}

\subsection{Scale}

Beyond the UNHCR figures cited in \Cref{sec:intro}, the population
relevant to the present argument is broader than the
forcibly-displaced and stateless categories alone. Undocumented
migrants who hold no recognized status in their country of
residence, internally displaced persons whose host
sub-jurisdiction does not extend banking infrastructure to their
location, and populations in failed-state or contested-sovereignty
zones whose nominal citizenship is not exercisable for commercial
purposes all share the operationally-relevant property: they
cannot reliably enter the legal infrastructure of commerce as
recognized parties. Estimates of the combined population vary
substantially with definition; the order of magnitude is
hundreds of millions, not millions.

\subsection{The Received Architecture}

Labor law, contract law, banking, and commercial-dispute
adjudication each presume a recognized legal subject attached to
a jurisdiction that enforces their rights. The relevant
populations occupy the negative space of that presumption.

\begin{itemize}
  \item The employment relation is typically unavailable to them
    formally: permission-to-work requirements gate access to the
    employee category, and undocumented work, where it occurs,
    occurs outside legal recognition.
  \item The independent-contractor relation is unstable: cannot
    register a business in many jurisdictions of residence, cannot
    enforce contracts in local courts that do not recognize the
    party's standing, cannot reliably serve process across
    jurisdictional boundaries.
  \item Banking access is restricted or absent: opening an account
    typically requires documents of residency and identity that
    the population does not have or does not have in the form the
    host jurisdiction recognizes.
  \item Adjudication has no stable forum: commercial disputes have
    no local court willing to hear them from a party whose
    standing is contested.
\end{itemize}

The humanitarian-economics tradition has documented the cost of
this gap. \citet{betts_collier_2017} argue that displaced
populations are systematically prevented from contributing to the
economies that host them by an institutional architecture
designed for citizens, with the result that the host economy
captures little of the productive capacity the displaced bring,
and the displaced themselves remain dependent on humanitarian
transfers when productive participation would be both possible
and preferred. \citet{pritchett2006people} makes the broader
argument about labor-mobility friction and the welfare cost of
the institutional gap between labor markets in origin and
destination economies.

\subsection{What the Primitive Does and Does Not Do}

The primitive does not resolve statelessness. It does not confer
citizenship, restore legal recognition, or create enforcement
against a state that denies the person's standing. It is not a
substitute for the political work that addresses those gaps.
What it does is structurally narrower:

\begin{enumerate}[label=(\alph*)]
  \item It provides a commerce architecture whose precondition is
    a cryptographic key rather than a recognized civil-legal
    subjecthood. A wallet does not require a passport to bond; a
    bonded commitment does not require a host-state court to
    enforce; a settled exchange does not require a domestic bank
    to clear.
  \item It produces a verifiable performance record that travels
    with the wallet across jurisdictional boundaries. A
    contributor whose civil-legal subjecthood is contested in
    jurisdiction $A$ accumulates the same on-chain attestation
    history as one whose status is settled in jurisdiction $B$.
    The record is portable in a way that conventional employment
    or contracting records are not.
  \item It opens a region of productive activity --- bilateral,
    bonded, settled directly between parties --- to participants
    for whom the legal scaffold that conventionally underwrites
    such activity is unavailable. The set of activities thus
    opened is bounded by what can be coordinated bilaterally and
    settled atomically; the bound is non-trivial but not all of
    productive activity.
\end{enumerate}

These are descriptive properties of the architecture. They do not
depend on any humanitarian program or any particular legal
reform; they follow from the construction of the primitive.

\subsection{Arendt and the Right to Have Rights}

\citet{arendt1951origins} observed that the fundamental human
right is the right to have rights --- the right to belong to a
political community within which one's other rights have
standing. Her analysis of statelessness named the condition as a
civilizational failure rather than an individual misfortune. The
person without a polity is not merely deprived of specific
rights; they are deprived of the institutional setting in which
having-rights is even a coherent claim. The remedy is political,
not technological.

The Arendtian objection to the present paper's reading is sharper
than the surface formulation, and the paper is improved by
engaging it directly. Part II Chapter 9 of \emph{The Origins of
Totalitarianism} (``The Decline of the Nation-State and the End
of the Rights of Man'') is not only a diagnosis of statelessness;
it is an analysis of a specific kind of \emph{worldlessness} --- the
loss not of particular entitlements but of the standing of the
human person within a public world that confers significance on
their actions. \citet{arendt1958human}'s later distinction
between \emph{labor}, \emph{work}, and \emph{action} sharpens
this: the world is constituted through action visible to others,
in a public space within which the actor is recognized as an
actor; loss of polity-membership is loss of the public space
within which action is constitutive of the world.

The Arendtian objection to a ``capacity to have commerce''
framing is therefore not that commerce is unimportant for the
displaced; it is that commerce-without-polity may be the
deepening of worldlessness rather than its alleviation. To
exchange goods with strangers across a substrate that confers no
membership and recognizes no public standing is, on a strict
Arendtian reading, to participate in \emph{labor} (the metabolic
necessity of staying alive) without re-entering \emph{action}
(the recognized public act). The bonded substrate looks, on this
reading, like a more efficient form of bare life rather than a
restoration of the world.

We take the objection seriously and decline it on these terms.
The capacity-to-have-commerce we describe is not a substitute
for membership; it is a structural precondition that membership
historically subsumed and that, where membership is denied, can
be supplied separately. The bonded transaction does include a
recognition function in Arendt's sense, on a small scale: the
commitment is dual-signed, the performance is attested, the
record is public, and the counterparty acknowledges the
participant as a party with whom binding agreements can be made.
This is not the political recognition Arendt regarded as
fundamental --- the recognition of co-membership in a polity ---
but it is not nothing in her register either. It is bilateral
recognition of capacity-to-bind, conferred not by a political
authority but by a counterparty who has chosen to bond with the
participant. The accumulated record of such recognitions across
many bonded relationships does not constitute a polity, but it
does constitute a public visibility of the participant's actions
that statelessness conventionally denies. The paper's claim is
therefore narrower than the Arendtian framing of restoration:
the substrate supplies a portable evidentiary subjecthood that
is not coextensive with political membership, but that fills a
gap statelessness has historically left empty.

The bonded primitive does not grant the right to have rights; it
grants, more narrowly, a capacity to have commerce. The two
should not be confused. The right to have rights is political and
is not within the gift of any technology. The capacity to have
commerce is structural and becomes operational once the
cryptographic infrastructure is in reach. For populations for
whom the first right is denied, the second capacity is
non-trivial; for populations for whom the first right is granted,
the second capacity is largely subsumed by existing institutions
and the primitive is one infrastructure among several. The
distinction matters for the audience of this paper because the
primitive's leverage is asymmetric: where the conventional
infrastructure works, it adds modestly; where the conventional
infrastructure fails, it adds the architectural condition for
participation at all.

\paragraph{The protection-gap residue the substrate cannot reach.}
The 1951 Refugee Convention and its 1967 Protocol enumerate a set
of rights that the host state owes to refugees within its
territory: \emph{non-refoulement} (Article~33), access to courts
(Article~16), identity papers (Article~27), travel documents
(Article~28), employment, public education, and public relief
\citep{refugee_convention_1951, goodwin_gill_mcadam_2021}. The
\emph{Right to Have Rights} as Arendt names it is closer to the
spirit of the Convention than to any single article; the
Convention operationalizes the spirit through specific institutional
duties \citep{kesby2012right, benhabib2004rights}. The bonded
primitive supplies none of these: it cannot prevent
\emph{refoulement} (which is a sovereign decision on territorial
admission); it cannot grant access to courts (which is a
jurisdictional decision); it cannot issue identity papers,
travel documents, or work permits. The protection gap that
\citeauthor{goodwin_gill_mcadam_2021} describe is real, vast, and
political; the substrate addresses one of its many components
(the inability to bind in commerce without a recognized
counterparty) and leaves the others where they were.

A reader sympathetic to refugee policy should take this honestly:
the present paper is not an argument for de-prioritizing the
political work of Convention enforcement, host-state obligations,
or international protection regimes. The substrate is a
\emph{useful adjunct} where it works, not a substitute for the
broader regime. The position we hope to occupy is closer to
\citeauthor{benhabib2004rights}: the rights-of-others framework
identifies layers of moral and legal claim that are differently
secured by different institutional arrangements, and the bonded
substrate adds one such layer (commerce-binding capacity) without
displacing the others.

% ════════════════════════════════════════════════════════════════════
\section{Operational Constraints}
\label{sec:constraints}

The present paper is positional rather than programmatic. The
primitive moves the commerce-access gap from architectural
impossibility to operational engineering question; it does not
solve the engineering question. We name four constraints that
remain.

\paragraph{Key custody.}
Populations with unreliable access to power, devices, and
connectivity face usability barriers at the cryptographic layer
that the primitive does not solve. Loss-of-key risk is high in
unstable environments, and key-recovery infrastructure for
displaced populations is not built. The constraint is real and
evolves with the broader infrastructure around the protocol;
solutions exist (social recovery, custodial wallets with
revocable trust) but each carries trade-offs the population is
poorly positioned to evaluate.

\paragraph{Counterparty willingness.}
A well-resourced buyer in a recognized jurisdiction may decline
to bond with a counterparty whose civil-legal standing is
contested. The decline is not because of the primitive; it is
because of downstream legal exposure in the buyer's home
jurisdiction (sanctions screening, anti-money-laundering
obligations, forum-of-enforcement concerns). The primitive does
not dissolve this friction. Its effect, where it has one, is to
shift the friction from architectural --- ``commerce with the
displaced is impossible'' --- to discretionary --- ``some
counterparties will, others will not.''

\paragraph{Fiat on-ramp.}
Converting wallet-denominated value to local currency still runs
through existing rails, which remain gated for the population in
question. The primitive operates within the crypto-native
settlement layer; the bridge between that layer and local
purchasing power remains a separate problem with its own
architecture and its own gatekeepers. Stablecoin acceptance by
local merchants, where it occurs, partially closes this gap; full
closure depends on infrastructure outside the primitive.

\paragraph{Collateral access.}
Bonding requires collateral. A counterparty without starting
capital cannot bond. This is a hard constraint on the primitive's
applicability to populations whose displacement is correlated with
asset loss. Possible mitigations exist at the assembly tier ---
sponsor pools, cooperative bonding, assembly-aggregated bonds,
microfinance-style group guarantees --- but each of these is a
construction on top of the primitive, not a property of the
primitive itself, and none are yet built at scale. This
constraint is not minor: the primitive's equilibrium results
assume both parties are legally and materially free to enter the
bonding
game --- that entering is a choice, and that capital sufficient
to post the required bond is available. For displaced populations
whose displacement is itself the loss of capital, the assumption
fails on the second clause. The primitive is symmetric in form
and asymmetric in the access conditions that make use of the
form possible. Among the four constraints listed, this is the
one most likely to determine whether the primitive's contribution
to the population in question is real or merely technical.

\paragraph{Bond-currency jurisdictional embedding.}
A constraint adjacent to the four above and more easily missed:
the bond's denomination currency is itself a jurisdictional
artifact. The primitive is denomination-agnostic, but the
denominations available in practice are not. The principal
US-dollar-pegged stablecoin (USDC) is issued by Circle, a
US-jurisdictional entity subject to OFAC sanctions enforcement
and capable of freezing addresses on Treasury request; the
\emph{Tornado Cash} sanctions of August 2022, which targeted a
piece of immutable open-source code rather than an entity, made
explicit how far the regulatory reach can extend
\citep{ofac2022tornado, vanvalkenburgh2022tornado}. ETH is a
neutral substrate at the protocol level but is acquired through
exchanges that KYC, and on-ramps from local fiat to ETH require
banking relationships that the displaced participant typically
lacks. A wallet that the bonded primitive cannot reach because
the currency-layer gatekeepers refuse to bridge to it is not
made reachable by anything the substrate does.

The choice of denomination therefore matters for the
displaced-participant case in a way it does not for the
counterparty case. We recommend, where the architecture admits
it, denominations whose issuance and transfer are not subject to
discretionary freeze authority (which excludes most fiat-pegged
stablecoins as currently constituted) and whose acquisition does
not run through KYC-gated exchanges (which excludes most retail
on-ramp paths). In practice the operational candidate is
permissionless on-chain assets with deep secondary-market
liquidity and no issuer-level freeze authority; the universe is
narrower than the universe of all crypto-assets, and the
selection criterion is structural rather than ideological.

The cumulative effect of these constraints is that the primitive
is not a solved-problem deliverable for displaced populations
today. Its present-tense contribution is architectural:
commerce-without-recognition is now possible to build, where it
was not before, and the constraints that remain are engineering
problems with engineering solutions rather than structural
features of the legal infrastructure.

\paragraph{Prior art the paper should engage.}
The constraints listed above are not theoretical: they have been
surfaced and partially addressed by humanitarian programs that
have piloted blockchain-based assistance over the past decade.
The UN High Commissioner for Refugees' \emph{Building Blocks}
program, deployed in Jordan since 2017 and expanded across
several refugee operations, uses on-chain transaction records to
coordinate cash assistance from multiple humanitarian agencies to
the same beneficiaries while reducing duplicate disbursement and
intermediary cost \citep{unhcr2024buildingblocks,
coppi_fast_2019}. The World Food Programme's similar pilots
have produced substantial operational experience with the
key-custody, fiat-on-ramp, and counterparty-recognition
constraints we name above. The independent evaluations
\citep{coppi_fast_2019, caribou_digital_2018} are mixed: the
programs achieve real operational improvements in coordination
and disbursement, but the binding constraint identified across
evaluations is not architectural recognition of the participant
on chain --- which the architecture supplies straightforwardly ---
but the host-state and on-ramp gatekeeping that we list above as
distinct constraints. The displaced beneficiaries in these
programs do not bond; they receive. The bonded primitive
extends the architecture to a use case the existing programs
have not yet exercised at scale: the displaced participant as
\emph{counterparty} in a productive bilateral exchange rather
than as recipient of one-sided disbursement. The empirical
record above is informative about what the engineering will
require; it is not yet informative about what will happen when
displaced participants are bonded counterparties, because the
relevant pilots have not been run.

% ════════════════════════════════════════════════════════════════════
\section{Conclusion}
\label{sec:conclusion}

The bonded settlement primitive does not solve statelessness, and
this paper has not argued that it does. It does change the
architecture of the commerce-access gap that statelessness
creates, by replacing the precondition of recognized
civil-legal subjecthood with the precondition of a cryptographic
key. The change is not in the political condition of the
populations under discussion; it is in the institutional
architecture they face when they attempt to participate in
productive activity across jurisdictional boundaries.

The Arendtian distinction is the one to hold: the right to have
rights is political, denied by polities, and not within the gift
of any technology to restore. The capacity to have commerce is
structural, denied by the legal architecture of commerce, and
restorable by a settlement primitive whose precondition is
cryptographic rather than civil-legal. For populations for whom
the first denial is constitutive, the second restoration is
non-trivial.

The paper is positional. It identifies a gap, names the
architectural change, and acknowledges the operational
constraints that remain. The programmatic work --- building
key-custody infrastructure for refugee populations, designing
sponsor pools that lend collateral against verified attestation
history, partnering with humanitarian organizations to evaluate
where bonded commerce is structurally feasible and where it is
not --- is for the practitioners who do that work. The
contribution of this paper is to make the
architectural change legible to scholarship that has not yet
recognized it as available.

\section*{Acknowledgments}
The argument here owes a particular debt to readers from
humanitarian-economics and refugee-policy backgrounds who
insisted the statelessness extension be treated separately from
the gig-economy classification debate, on the grounds that
conflating the two did neither audience justice. The Arendtian framing came from
conversations with collaborators working on the legal-philosophy
side of statelessness scholarship; specific responsibility for the
present paper's positioning rests with the author.


% ════════════════════════════════════════════════════════════════════
% References
% ════════════════════════════════════════════════════════════════════

\begin{thebibliography}{99}

\bibitem[Arendt(1951)]{arendt1951origins}
Arendt, H.
\newblock \emph{The Origins of Totalitarianism}.
\newblock Harcourt, Brace \& Co., New York, 1951.
\newblock See especially Part II, Ch.~9 (``The Decline of the
  Nation-State and the End of the Rights of Man''), for the
  formulation of the right to have rights.

\bibitem[Betts and Collier(2017)]{betts_collier_2017}
Betts, A. and Collier, P.
\newblock \emph{Refuge: Transforming a Broken Refugee System}.
\newblock Penguin, London, 2017.

\bibitem[Pritchett(2006)]{pritchett2006people}
Pritchett, L.
\newblock \emph{Let Their People Come: Breaking the Gridlock on
  Global Labor Mobility}.
\newblock Center for Global Development, Washington, DC, 2006.

\bibitem[Allen(2016)]{allen2016path}
Allen, C.
\newblock The Path to Self-Sovereign Identity.
\newblock Personal blog, April 25, 2016.
\newblock \url{https://www.lifewithalacrity.com/article/the-path-to-self-sovereign-identity/}

\bibitem[Buterin(2021)]{buterin2021socialrecovery}
Buterin, V.
\newblock Why We Need Wide Adoption of Social Recovery Wallets.
\newblock Personal blog, January 11, 2021.

\bibitem[Tobin and Reed(2017)]{tobin_reed_2017}
Tobin, A. and Reed, D.
\newblock The Inevitable Rise of Self-Sovereign Identity.
\newblock Sovrin Foundation White Paper, 2017.

\bibitem[W3C(2022)]{w3c_did_2022}
W3C.
\newblock Decentralized Identifiers (DIDs) v1.0: Core Architecture,
  Data Model, and Representations.
\newblock W3C Recommendation, July 2022.

\bibitem[Locke(1689)]{locke1689treatises}
Locke, J.
\newblock \emph{Two Treatises of Government}.
\newblock Awnsham Churchill, London, 1689.

\bibitem[Nozick(1974)]{nozick1974anarchy}
Nozick, R.
\newblock \emph{Anarchy, State, and Utopia}.
\newblock Basic Books, New York, 1974.

\bibitem[Argent(2020)]{argent_social_recovery}
Argent.
\newblock Social Recovery: Reinventing Wallet Security.
\newblock Argent Documentation, ongoing.

\bibitem[Safe(2022)]{safe_social_recovery}
Safe.
\newblock Social Recovery via Threshold Multi-Signature.
\newblock Safe (formerly Gnosis Safe) Documentation, ongoing.

\bibitem[Arendt(1958)]{arendt1958human}
Arendt, H.
\newblock \emph{The Human Condition}.
\newblock University of Chicago Press, Chicago, 1958.

\bibitem[Benhabib(2004)]{benhabib2004rights}
Benhabib, S.
\newblock \emph{The Rights of Others: Aliens, Residents, and
  Citizens}.
\newblock Cambridge University Press, Cambridge, 2004.

\bibitem[Goodwin-Gill and McAdam(2021)]{goodwin_gill_mcadam_2021}
Goodwin-Gill, G.~S. and McAdam, J.
\newblock \emph{The Refugee in International Law}, 4th edition.
\newblock Oxford University Press, Oxford, 2021.

\bibitem[Kesby(2012)]{kesby2012right}
Kesby, A.
\newblock \emph{The Right to Have Rights: Citizenship, Humanity, and
  International Law}.
\newblock Oxford University Press, Oxford, 2012.

\bibitem[Refugee Convention(1951)]{refugee_convention_1951}
United Nations.
\newblock \emph{Convention Relating to the Status of Refugees}.
\newblock 189 UNTS 137, 1951; expanded by the 1967 Protocol,
  606 UNTS 267.

\bibitem[OFAC(2022)]{ofac2022tornado}
US Department of the Treasury, Office of Foreign Assets Control.
\newblock \emph{Treasury Sanctions Notorious Virtual Currency
  Mixer Tornado Cash}.
\newblock Press release, August~8, 2022.

\bibitem[Van Valkenburgh(2022)]{vanvalkenburgh2022tornado}
Van Valkenburgh, P.
\newblock Analysis of the Tornado Cash Sanctions: How Far Does
  Treasury's Authority Reach?
\newblock Coin Center Report, 2022.

\bibitem[UNHCR(2017)]{unhcr2024buildingblocks}
UN High Commissioner for Refugees.
\newblock \emph{Building Blocks: Blockchain for Humanitarian
  Assistance}.
\newblock UNHCR Innovation Service, ongoing since 2017.

\bibitem[Coppi and Fast(2019)]{coppi_fast_2019}
Coppi, G. and Fast, L.
\newblock Blockchain and Distributed Ledger Technologies in the
  Humanitarian Sector.
\newblock Humanitarian Policy Group Commissioned Report, Overseas
  Development Institute, 2019.

\bibitem[Caribou Digital(2018)]{caribou_digital_2018}
Caribou Digital.
\newblock \emph{Blockchain for Development: Hope or Hype?}
\newblock Caribou Digital Publishing, 2018.

\bibitem[UNHCR(2024)]{unhcr2024globaltrends}
United Nations High Commissioner for Refugees.
\newblock \emph{Global Trends: Forced Displacement in 2023}.
\newblock UNHCR, Geneva, 2024.

\end{thebibliography}

\end{document}
